\scanpage{9}
% archive foliation: 8
$r$ de $X$ sur $S$ est $1$, \deleted{\unsure{on a}} (alors $v_1,-v_1$ des iso) ou $r+1$ inv. sur $S$ (alors $v_1,-v_1$ ont des inv. à g.), \deleted{le} \inserted{$\Omega$ et $\Omega$-}torseur $\mathcal{T}$ est canoniquement déterminé par $X$, par $\Omega\simeq\Omega^1_{X/S}$ et $\unsure{\alpha_n}$.

On définit de même une $\infty$-connexion \inserted{projective} (Remarque : est-il vrai que pour une telle $\infty$-conn. proj, $\mathcal{T}$ est can. déterminé par $X/S$, sans condit. de dim ni de car. ?)

Notons que \inserted{la pro-$\mathcal{O}_X$-alg} $P^\infty_{X/S}$ a une \deleted{connexion} \inserted{stratification} \deleted{plate} canonique $\Delta$ (\deleted{\unsure{donc induit une connexion}} \inserted{ne respectant pas l'}augmentation $P^\infty_{X/S}\to\unsure{\mathcal{O}_X|\mathcal{O}_{X/S}}$ !), donc si $X$ est muni d'une $\infty$-connexion projective, alors \inserted{la pro-$\mathcal{O}_X$-alg} $\widehat P$ i.e. $\widehat S$ se trouve sur $X$ muni d'une \deleted{$\infty$-connexion} \inserted{stratification} \deleted{correspondante} \inserted{$\Delta'$}. On dit que la $\infty$-conn. projective est \inserted{plate (ou} \uline{intégrable}) si $\Delta'$ provient d'une \deleted{$\infty$-connexion} stratification de $P=\mathbb{P}(E)$ (en oubliant \unsure{pour} la section marquée à présent) laquelle sera dès lors uniquement déterminée (et qualifiée de ``canonique'').

Cette condition étant locale sur $X$, pour l'exprimer on peut supposer que $X$ est muni de ``coordonnées locales'' $f_1\ldots f_r$ (\unsure{sections relat. à un $\alpha$ sur $S$}). Dès lors
\[
P^\infty_{X/S}\simeq\mathcal{O}_S[[\delta f_1,\ldots,\delta f_r]]
\]
et
\[
\Delta:P^\infty_{X/S}\longrightarrow P^\infty_{X/S}\widehat\otimes_{\mathcal{O}_S}P^\infty_{X/S}
\simeq P^\infty_{X/S}(P^\infty_{X/S})\simeq P^\infty_{X/S}[[\overline{\delta f_1},\ldots,\overline{\delta f_r}]]
\]
\[
\Delta(\delta f_i)=\delta(f_i)\otimes1+1\otimes\delta(f_i)=\delta f_i+\overline{\delta f_i}
\]
\margin{\[
M\otimes P^n_{X/S}\longrightarrow P^n_{X/S}\otimes M
\]
\deleted{un} $P^n_{X/S}$-linéaire qui est \deleted{\unsure{bij}} \inserted{par réduction}
i.e.
\[
M\xrightarrow{\Delta}P^n_{X/S}\otimes M
\]
$\mathcal{O}_S$-lin par stratification. id. après réduction \inserted{sur 2ème facteur}
\[
\Delta(m)=m+\sum
\]}

On peut supposer l'extension $E$ splittée, donc $\widehat S\simeq\widehat{\operatorname{Sym}}(\Omega)$ et comme $\Omega\simeq\Omega^1_{X/S}\simeq\mathcal{O}_S^r$,
\[
\widehat S\simeq\mathcal{O}_S[[T_1,\ldots,T_r]]
\]
La $\infty$-conn. projective est
